\documentclass[output=paper]{langscibook}
\ChapterDOI{10.5281/zenodo.22880286}
\title{\textit{Kakavei} – The bamboo frame net}
\author{Naphtali Maion}
\abstract{}
\begin{document}
\maketitle\label{chap:29}
\begin{figure}[H]
    \includegraphics[trim=0 0 0 0,clip,width=.9\linewidth]{figures/29_1.jpg}
    \caption{\textit{Kakavei} -- Bamboo-frame net}
    \label{fig:29_1}
\end{figure}
\begin{figure}[H]
    \includegraphics[trim=0 0 0 0,clip,width=.9\linewidth]{figures/29_2.png}
    \caption{\textit{Hihiana} -- Net needle}
    \label{fig:29_2}
\end{figure}
\begin{xltabular}{\textwidth}{lXX}
    (\ref{ex:text22-1})&\textbf{\textit{A kakavei}}&\textbf{The \textit{kakavei} net}\\
	(\ref{ex:text22-2})&\textit{Vaamatatopo na maa taba}&The preparation of the things (you need to make a pair of \textit{kakavei} nets):\\
	(\ref{ex:text22-3})&\textit{Momohu kona nia kave tabae tau paku an, a kave a beera ge a rutaa}&First, think about what kind of net you are going to make, a large net or a small one.\\
	(\ref{ex:text22-4})&\textit{Beo hum bona kavara, erau ean antee vai tea kaokao a maamihu taba na kakavei tean.}&When this part is finished, then you can go ahead with all the things of your \textit{kakavei}.\\
	(\ref{ex:text22-5})&\textit{Kona ni amaa humhum na kakavei.}&Think of the parts of the \textit{kakavei} net.\\
	(\ref{ex:text22-6})&\textit{A hena na maa taba paku na kakavei:}&The names of the things for making a \textit{kakavei} net:\\
	(\ref{ex:text22-7})&\textit{1. koverau (baoru) -- 2}&two bent \textit{koverau} bamboos\footnote{See Chapter \ref{chap:10}:\ref{ex:text6-1}ff.}\\
	(\ref{ex:text22-8})&\textit{2. kapikapis -- 8}&eight wedges (made of wood)\\
	(\ref{ex:text22-9})&\textit{3. kara raviana -- 6}&six strings made of \textit{raviana} vines\footnote{See Chapter \ref{chap:11}:\ref{ex:text7-75}ff.}\\
    (\ref{ex:text22-10})&\textit{4. toorana (naono) -- 2}&two \textit{toorana} poles to which the bamboo is tied\\
	(\ref{ex:text22-11})&\textit{5. sunsun (miimiana) -- 2}&two \textit{sunsun} poles of \textit{miimiana} palm wood\footnote{See \url{https://dictionaria.clld.org/units/teop-miimiana}},\\
	(\ref{ex:text22-12})&\textit{6. kara baneo ge kara kave -- bero}&many strings made of fibres of the young breadfruit  tree or the net vine\footnote{See Chapter \ref{chap:8}:\ref{ex:text4-1}ff, Chapter \ref{chap:11}:\ref{ex:text11-17}ff, Chapter \ref{chap:27}:\ref{ex:text20-3}ff.}\\
\end{xltabular}	
\begin{xltabular}{\textwidth}{lXX}
	(\ref{ex:text22-13})&\textit{7. tavara (kururu) -- 1}&a small rectangular piece of \textit{kururu} bamboo (2 x 4 inches) that is used to measure the size of the meshes of the net\\
    (\ref{ex:text22-14})&\textit{8. hihiana (o koverau o rutaa)}&a net needle (made of a small piece of \textit{koverau} bamboo)\\
	(\ref{ex:text22-15})&\textit{Bean tau paku a kakavei,}&When you are about to make a \textit{kakavei},\\
	(\ref{ex:text22-16})&\textit{evei amaa taba toro vaamatatopo an.}&these are the things that you must prepare.\\
	(\ref{ex:text22-17})&\textit{(1) O buaku o koverau mene o buo baoru.}&Two \textit{koverau} bamboos for two bamboo frames.\\
	(\ref{ex:text22-18})&\textit{Mosi amaa koverau vai amaa manoto vamataa.}&Cut fully matured \textit{koverau} bamboos.\\
	(\ref{ex:text22-19})&\textit{Vaaraka vamatai tea roava.}&Let them thoroughly dry in the sun.\\
	(\ref{ex:text22-20})&\textit{Taraha taatagi amaa koverau bona.}&Carefully level the outside of these \textit{koverau} bamboos (i.e. the knots).\\
	(\ref{ex:text22-21})&\textit{(2) O buaku o toorana toro mosi maa. Vaarakai tea roava.}&The two \textit{toorana} poles must be cut. Let them dry in the sun.\\
	(\ref{ex:text22-22})&\textit{(3) Mosi maa o miimiana mene o buaku o sunsun.}&Cut a \textit{miimiana} palm for the two \textit{sunsun} poles.\\
	(\ref{ex:text22-23})&\textit{Vaarakai mee tea roava.}&Let it also dry in the sun.\\
	(\ref{ex:text22-24})&\textit{(4) Mosi a vahara baneo, kahu a kapa nae, vaatei tea roava ore raka.}&Cut young breadfruit trees, scrape off the bark, put them into the sun to dry.\\
	(\ref{ex:text22-25})&\textit{Be raka, hihi a kara nae mene paku bona kanono.}&When they are dry, pull their fibres off to make the string.\\
	(\ref{ex:text22-26})&\textit{Bean rake, ean repaa kurus kara kave maa rasuu.}&If you want, you cut net vines in the jungle.\\
	(\ref{ex:text22-27})&\textit{Vaarakai tea roava, ean repaa kahue bona kehaa.}&Let them dry in the sun, and then you scrape them with a seashell.\\
	(\ref{ex:text22-28})&\textit{Ean repaa paku a kanono mene huha bona kakavei.}&Then you make the string for tying the \textit{kakavei} net.\\
	(\ref{ex:text22-29})&\textit{(5) Tasu varaba maa a raviana rasuu.}&Cut the \textit{raviana} vines in advance in the jungle.\\
	(\ref{ex:text22-30})&\textit{Vaarakai tea roava bean siri kopaa varaba batae.}&Let them dry in the sun after you have torn them lengthwise apart.\\
	(\ref{ex:text22-31})&\textit{(6) O kapikapis: Amaa ta naono vaarii na taraha vakavara raara mene kapis bona maa pampam na kakavei.}&The wedges: we carve all these pieces of wood to wedge them into the lashings of the \textit{kakavei} net.\\
	(\ref{ex:text22-32})&\textit{Taraha vaantee eta kukan totoka to antee nana bona buaku a kakavei.}&Carve eight pieces which will be enough for the two \textit{kakavei} nets.\\
	(\ref{ex:text22-33})&\textit{(7) Vaamatatopo varaba a tavara.}&Prepare the \textit{tavara} (piece of wood for measuring the meshes of a net).\\
	(\ref{ex:text22-34})&\textit{(8) Vaamatatopo varaba o koverau o rutaa mene o hihiana.}&Prepare a small piece of \textit{koverau} bamboo in advance for the net needle.\\
\end{xltabular}
\begin{xltabular}{\textwidth}{lXX}
	(\ref{ex:text22-35})&\textit{\textbf{Vaamatatopo a kanono}}&\textbf{Prepare the string}\\
	(\ref{ex:text22-36})&\textit{O hum kiu vai na beera oha vareko nana.}&This part of the work is quite labour intensive.\\
	(\ref{ex:text22-37})&\textit{Ean kahi tagihu nom rasuu tea kurus kave.}&You go into the jungle to cut net-vines.\\
	(\ref{ex:text22-38})&\textit{O kave na abeabe batana tea naono, saka tei vananaoana vira haana me.}&The net vine is creeping on the trees, it is also not very far away.\\
	(\ref{ex:text22-39})&\textit{Bean suguna koa to abe vaberobero riori,}&When you have arrived where they are creeping in large quantities,\\
	(\ref{ex:text22-40})&\textit{ean repaa saku koa tea kurus, ean repaa voosu vahaa koa.}&you just quickly cut them, and then you go back home.\\
	(\ref{ex:text22-41})&\textit{Bean voosu me amaa dee kave, pasi vaobete roho noman,}&When you come home with loads of net-vines, you first lay them down,\\
	(\ref{ex:text22-42})&\textit{eve re meapa roho.}&so that they become flexible.\\
	(\ref{ex:text22-43})&\textit{Beve meapa, ean repaa hihi manae,}&When they are flexible, you eventually remove the fibres from them,\\
	(\ref{ex:text22-44})&\textit{ean repaa kahu batae mee.}&then you also scrape them.\\
	(\ref{ex:text22-45})&\textit{Bean kahu vakavarae,}&When you have completely scraped them,\\
	(\ref{ex:text22-46})&\textit{ean repaa varoava e ore raka.}&you put them into the sun to dry.\\
	(\ref{ex:text22-47})&\textit{Be raka vamataa, ean repaa viru a kanono tean.}&When they have dried up well, you twist your string.\\
	(\ref{ex:text22-48})&\textit{A kanono bona to kahi huha noman bona kakavei tean.}&This string (is the one) that you are going to tie your \textit{kakavei} net with.\\
	(\ref{ex:text22-49})&\textit{A kanono toro antee bona buaku a kakavei a bua vapara.}&The string must be (long) enough for the pair of \textit{kakavei}-nets.\\
\end{xltabular}
\begin{xltabular}{\textwidth}{lXX}
	(\ref{ex:text22-50})&\textbf{\textit{Huha kakavei}}&\textbf{Tying \textit{kakavei}-nets}\\
	(\ref{ex:text22-51})&\textit{Bea maa taba vai matatopo, ean antee vai tea huha a kave.}&When these things are prepared, you can now tie the net.\\
	(\ref{ex:text22-52})&\textit{A taba to pasi paku noman, ei a kanono vai to kahi huha noman bona buaku a kave.}&The thing that you will make, this is the string that you use for tying the two nets.\\
	(\ref{ex:text22-53})&\textit{A kanono bona pasi vaabebeahu kurus noman ore antee bona buaku a kave.}&You will make this string very long so that it is long enough for the two nets.\\
	(\ref{ex:text22-54})&\textit{A meha kanono to kahi paku noman, ei a roveana.}&Another string that you'll make is the \textit{roveana}.\\
	(\ref{ex:text22-55})&\textit{A roveana eve a kanono to taviri vakavara nana bona kave.}&The \textit{roveana} is the string that surrounds the whole net as a seam.\footnote{Compare Chapter \ref{chap:27}, \figref{fig:27_1}.}\\
	(\ref{ex:text22-56})&\textit{A kanono bona ge a roveana to rahi parata nana bona kave teo vuri to nove noman a kakavei.}&This string or \textit{roveana} pulls the net open when you tie the \textit{kakavei} net to the frame.\\
	(\ref{ex:text22-57})&\textit{Bea bua kanono vai kavara, ean matatopo vai tea paku a sinoo na kave.}&When the two strings are finished, you are prepared to make the frame of the net.\\
	(\ref{ex:text22-58})&\textit{A kanono tean a matatopo,}&Your string is prepared,\\
	(\ref{ex:text22-59})&\textit{ean viru vai a kanono tea hihiana.}&you  now wind the string now around the net needle.\\
	(\ref{ex:text22-60})&\textit{Ean gono vai a tavara.}&You get the \textit{tavara} (to measure the size of the meshes).\\
\end{xltabular}
\begin{xltabular}{\textwidth}{lXX}
	(\ref{ex:text22-61})&\textit{\textbf{A sinoo na kave:} vatotomo na maa humhum na kakavei}& \textbf{The frame of the net:} putting the parts of the \textit{kakavei} net together\\
	(\ref{ex:text22-62})&\textit{A maamihu taba paa vaamatatopo vaevuru, a kakavei kahi vatotomo raarai.}&Everything is aready prepared, we will put the \textit{kakavei} net together now.\\
	(\ref{ex:text22-63})&\textit{(1) Gono o toorana bara o sunsun, pampam vabuaku ri bona kara raviana.}&Get the \textit{toorana} pole and the sunsun pole, tie them together with two lashings of the \textit{raviana} strings.\\
	(\ref{ex:text22-64})&\textit{(2) Gono o baoru, pam mi nai bono pasu no toorana.}&Get a bamboo pole, lash it together with the end of the \textit{torana}-pole.\\
	(\ref{ex:text22-65})&\textit{Pam i bona kara raviana (a buaku a pam).}&Lash them together with the \textit{raviana} strings (two lashings).\\
	(\ref{ex:text22-66})&\textit{(3) Tasu vakikis amaa pampam bona maa kapikapis.}&Firmly tighten the lashings with the wedges.\\
	(\ref{ex:text22-67})&\textit{Teo peho sinoo tapa ni o dorana o hum.}&Drive four pieces (i.e. wedges) into the lashings of one frame.\\
	(\ref{ex:text22-68})&\textit{Be kavara, o sinoo paa matatopo vaevuru.}&When it is finished, the frame is already prepared.\\
	(\ref{ex:text22-69})&\textit{(4) Nove a kakavei.}&Tie the \textit{kakavei} net to the frame.\\
\end{xltabular}
\begin{xltabular}{\textwidth}{lXX}
	(\ref{ex:text22-70})&\textit{\textbf{Vuri nove na kakavei}}&\textbf{Steps of tying the \textit{kakavei} net to its frame}\\
	(\ref{ex:text22-71})&\textit{(1) Vakisi nio pasu kave teo pasu no baoru.}&Tie an edge of the net to the end of the bamboo pole.\\
	(\ref{ex:text22-72})&\textit{(2) Vakisi nio meho teo sunsun.}&Tie another (edge) to the \textit{sunsun} pole.\\
	(\ref{ex:text22-73})&\textit{(3) Vakisi nio meho pasu kave romuno baoru.}&Tie the third edge of the net to the middle of the bamboo pole.\\
	(\ref{ex:text22-74})&\textit{(4) Vakisi nio vakosina teo toorana.}&Tie the end of net to the \textit{toorana} pole.\\
	(\ref{ex:text22-75})&\textit{A kakavei paa taparata vamataa vaevuru.}&The \textit{kakavei} has now already properly unfolded.\\
	(\ref{ex:text22-76})&\textit{Bean toku, tara teo vaa'abeabe na kakavei gina.}&If you do not understand it, look at the picture of the \textit{kakavei} above.\\
\end{xltabular}
\end{document}